\chapter{Metode Penelitian}
\emph{Cara menulis bab ini.} Bab Metode Penelitian menjelaskan bagaimana penelitian dilakukan: data, variabel, alur/tahapan, serta teknik analisis. Uraikan cukup rinci agar penelitian dapat direplikasi.


\section{Kerangka Pikir}
\emph{Cara menulis.} Gambarkan keterkaitan Permasalahan, Tujuan Penelitian, Solusi yang Diajukan, dan Kriteria Evaluasi; sajikan sebagai diagram yang dibuat ulang sendiri (tidak menyalin langsung dari sumber lain).

Kerangka pikir penelitian disajikan pada Gambar~\ref{gbr:kerangka}.

\begin{figure}[H]
\centering
\begin{tikzpicture}[node distance=1.1cm]
  \node[fc-mulai] (a) {Permasalahan};
  \node[fc-proses, below=of a] (b) {Tujuan Penelitian};
  \node[fc-proses, below=of b] (c) {Solusi yang Diajukan};
  \node[fc-io, below=of c] (d) {Kriteria Evaluasi};
  \draw[fc-panah] (a)--(b); \draw[fc-panah] (b)--(c); \draw[fc-panah] (c)--(d);
\end{tikzpicture}
\caption{Kerangka pikir penelitian}
\label{gbr:kerangka}
\end{figure}

\section{Data dan Variabel}
\emph{Cara menulis.} Gambarkan sumber dan cakupan data (mis.\ Susenas/Sakernas/PODES) serta variabel yang digunakan beserta definisinya.

Data utama bersumber dari Survei Sosial Ekonomi Nasional (Susenas) yang diperoleh
dari Badan Pusat Statistik (BPS). Variabel respons adalah pengeluaran per kapita.

\section{Metrik Evaluasi}
\emph{Cara menulis.} Jelaskan metrik evaluasi yang dipakai beserta rumusnya (mis.\ akurasi, \emph{precision}, \emph{recall}, F1).

Kinerja klasifikasi dievaluasi menggunakan akurasi. Sebagai contoh persamaan di
tengah paragraf, akurasi dihitung sebagai proporsi prediksi yang benar
$\mathrm{akurasi} = \frac{TP+TN}{P+N}$ dengan $TP,TN,P,N$ berturut-turut menyatakan
banyaknya prediksi positif benar, negatif benar, total observasi positif, dan
total observasi negatif.

Sebagai contoh persamaan pada baris tersendiri (perhatikan jaraknya yang rapat):
\begin{equation}\label{eq:f1}
F_1 = 2 \times \frac{\mathrm{precision}\times \mathrm{recall}}{\mathrm{precision}+\mathrm{recall}}.
\end{equation}
Persamaan~\ref{eq:f1} menjadi dasar evaluasi model. Untuk persamaan yang terlalu
panjang, operator penghubung diletakkan di AWAL baris lanjutan:
\begin{equation}\label{eq:model}
\begin{aligned}
\hat{y} ={}& \beta_0 + \beta_1 x_1 + \beta_2 x_2 + \beta_3 x_3 + \beta_4 x_4 \\
         &{}+ \beta_5 x_5 + \beta_6 x_6 .
\end{aligned}
\end{equation}

\section{Prosedur Estimasi}
\emph{Cara menulis.} Jelaskan teknik estimasi/algoritma yang dipakai, asumsi, dan kriteria pemilihan model (mis.\ metrik evaluasi, validasi silang).

Keseluruhan prosedur diringkas pada Algoritme~\ref{alg:pmt}.

\begin{algorithm}[H]
\caption{Estimasi dan evaluasi model PMT}\label{alg:pmt}
\KwIn{Data rumah tangga $\mathbf{X}$, target $y$, jumlah lipatan $K$}
\KwOut{Model terbaik $\hat{f}^\star$ dan metrik evaluasi}
Praproses data (imputasi, penskalaan, penyandian)\;
Bagi data menjadi himpunan latih dan uji\;
\ForEach{algoritma kandidat $m$}{
  \For{$k \leftarrow 1$ \KwTo $K$}{
    Latih $m$ pada $K-1$ lipatan\;
    Hitung metrik pada lipatan ke-$k$\;
  }
  Rerata metrik validasi silang untuk $m$\;
}
Pilih model dengan $F_1$ tertinggi sebagai $\hat{f}^\star$\;
\KwRet $\hat{f}^\star$\;
\end{algorithm}
